% =============================================================================
% Chapter 9  Discussion  --  ~4.5 pp
% rubric: Project Body (Modelling) -- Results & Discussion /30: "insight into
% the significance of the work", the band-4 separator.
% Plan: ThesisPlanning.md §5.9.
%
% PURPOSE     Explain and judge; do not restate Chapter 8.
% SOURCES     Planning.md §3.9 (complexity picture), §13 (risks); Chapter 8.
% WRITE WHEN  Gate G6 (wave W5).
% =============================================================================

\subsection{Why the Descriptor Form Wins, and What It Costs}
\label{sec:disc-tradeoff}
% ~1 pp
% LAND: norm, not encoding -- eliminating the derivatives multiplies norms
%   (||G|| = Theta(N^2), squared by the Gram product and again per derivative),
%   while carrying them keeps each block O(N) in norm. The trade: Hilbert space,
%   a block register, and rescaling where derivative norms dominate.

\subsection{End-to-End Complexity}
\label{sec:disc-complexity}
% ~1 pp
% LAND, in order:
%   1. The corrected complexity argument IN YOUR OWN VOICE: d = Theta(sqrt(
%      alpha_H / Delta) log 1/eps) (eq:degree), and alpha_H >= ||H|| for ANY
%      block encoding of this H.
%   2. So, for THIS FILTER FAMILY, the standard form cannot go below
%      Theta~(N^4) (eq:standard-wall) and the descriptor reaches Theta~(N)
%      (eq:descriptor-degree). Total gates Theta~(N poly n) against
%      Theta~(N^4 poly n).
%   3. Neither is poly-log. Say it first, yourself.
%   4. Hedge the "for any filter" generalisation to what Lin & Tong's lower
%      bounds actually support (lintong2020groundstate; ThesisPlanning.md §12,
%      risk TR-3). Consistent with sec:pihm-assessment, sec:standard-hamiltonian
%      and sec:descriptor-spectrum.

\subsection{What the Reproduction Says About the Published Method}
\label{sec:disc-reproduction}
% ~0.75 pp. Collegial and evidential register.
% LAND: a fair verdict -- the conventions are sound and the eta values right;
%   the printed choices distort at the printed n; the PDE and NDE panels do not
%   reproduce; the normalisation is loose; the readout identity needs a
%   correction for probabilistic preparation. Then what that means for readers
%   of the paper.

\subsection{Where the Construction Applies}
\label{sec:disc-applicability}
% ~0.5 pp
% LAND: guidance a practitioner could act on -- which equations benefit; when
%   block rescaling is needed; the nonlinear scope (dissipative, rho < 1).

\subsection{Feasibility on Fault-Tolerant Hardware}
\label{sec:disc-feasibility}
% ~0.5 pp
% LAND: logical qubits, CNOT or T counts and depth at meaningful n, from the
%   resource estimates (T3); why NISQ is out of scope.

\subsection{Limitations}
\label{sec:limitations}
% ~0.75 pp. Each limitation owned, WITH ITS ANALYSIS (the rubric's condition for
% no penalty):
%   no hardware; circuit simulation (T1) only at small n; the known-gap
%   assumption; the DCT as a cited primitive; the doubled space projects; NS
%   bounded by T3 at large K; the descriptor's Hilbert-space cost; Figs 6-7
%   unreproduced (authors not contacted).

% CHECKLIST (marker)
%   1. Does it explain rather than restate Chapter 8?
%   2. Is the complexity argument complete, and consistent with
%      sec:pihm-assessment, sec:standard-hamiltonian and sec:descriptor-spectrum?
%   3. Is every limitation owned, with its analysis?
%   4. Are consequences drawn for the field, not just for this project?
%   5. Is there actionable guidance a practitioner could follow?
